\documentclass[11pt,twoside,twocolumn]{usgsreport}
\usepackage{usgsfonts}
\usepackage{usgsgeo}
\usepackage{usgsidx}
\usepackage[figuretoc,tabletoc]{usgsreporta}

\usepackage{amsmath}
\usepackage{bm}
\usepackage{calc}
\usepackage{natbib}
\usepackage{graphicx}
\usepackage{longtable}

%Do not allow a page break to result in a line appearing by itself
% https://tex.stackexchange.com/questions/4152/how-do-i-prevent-widow-orphan-lines
\usepackage[all]{nowidow}

\makeindex
\usepackage{setspace}
\usepackage{etoolbox}
\usepackage{verbatim}

% set up the listings package for highlighting input files
\usepackage{listings}
\usepackage{xcolor}
\lstset{
  basicstyle=\footnotesize\ttfamily\color{black},
  numbers=none,
  columns=flexible,
  backgroundcolor=\color{yellow!10},
}
\lstdefinestyle{inputfile}{
  morecomment=[l]\#,
  backgroundcolor=\color{gray!10},
}
\lstdefinestyle{pythoncode}{
  language=Python,
  basicstyle=\scriptsize\ttfamily,
  keywordstyle=\color{blue},
  commentstyle=\color{gray!60!black},
  stringstyle=\color{red!60!black},
  showstringspaces=false,
  breaklines=true,
  backgroundcolor=\color{gray!8},
}
\lstdefinestyle{pseudocode}{
  basicstyle=\footnotesize\ttfamily,
  commentstyle=\color{gray!60!black},
  showstringspaces=false,
  breaklines=true,
  columns=fullflexible,
  backgroundcolor=\color{blue!5},
}

\usepackage{hyperref}
\hypersetup{
    pdftitle={MODFLOW 6 -- Description of Variables Modifiable Through the API},
    pdfauthor={MODFLOW 6 Development Team},
    pdfsubject={numerical simulation groundwater flow},
    pdfkeywords={groundwater, MODFLOW, simulation, API, BMI, XMI},
    pdflang={en-US},
    bookmarksnumbered=true,
    bookmarksopen=true,
    bookmarksopenlevel=1,
    colorlinks=true,
    allcolors={blue},
    pdfstartview=Fit,
    pdfpagemode=UseOutlines,
    pdfpagelayout=TwoPageRight
}

% reuse the cover image and figures from the mf6io guide
\graphicspath{{../mf6io/Figures/}}
\input{../version.tex}

\newcommand{\mli}[1]{\mathit{#1}}
\renewcommand\labelitemi{\tiny$\bullet$}

\renewcommand{\cooperator}
{Water Availability and Use Science Program}

\renewcommand{\reporttitle}
{MODFLOW 6 -- Description of Variables Modifiable Through the Application Programming Interface (API)}

\renewcommand{\preface}
{This document describes which internal MODFLOW 6 variables can be meaningfully modified at run time through the MODFLOW 6 Application Programming Interface (API), on a package-by-package basis. The API exposes MODFLOW 6 memory through the Basic Model Interface (BMI) and its eXtended Model Interface (XMI), which are most commonly accessed from Python through the \href{https://github.com/MODFLOW-ORG/modflowapi}{modflowapi} and \href{https://github.com/Deltares/xmipy}{xmipy} packages. The API exposes essentially all variables stored in the MODFLOW 6 memory manager, but not every variable can be changed with the expected effect: some are derived quantities that are computed once and then used directly, some are recomputed every time step and silently overwrite any value set externally, and some are re-read from input at the start of each stress period. This document classifies the variables in each package so that users of the API know which variable to set, and when, to achieve a desired change. The performance of MODFLOW 6 has been tested in a variety of applications. Future applications, however, might reveal errors that were not detected in the test simulations. Users are requested to send notification of any errors found in this report or in the model program to the contact listed on the \href{https://doi.org/10.5066/P1IJAXDZ}{MODFLOW Web page}. Additional details and information about ongoing developments are available through the version-controlled \href{https://github.com/MODFLOW-ORG/modflow6}{MODFLOW 6 repository}.}

\renewcommand{\coverphoto}{coverimage.jpg}
\renewcommand{\GSphotocredit}{Binary computer code illustration.}
\renewcommand{\reportseries}{}
\renewcommand{\reportnumber}{}
\renewcommand{\reportyear}{2017}
\ifdef{\reportversion}{\renewcommand{\reportversion}{\currentmodflowversion}}{}
\renewcommand{\theauthors}{MODFLOW 6 Development Team}
\renewcommand{\thetitlepageauthors}{\theauthors}
\renewcommand{\reportcitingtheauthors}{\theauthors}
\renewcommand{\colophonmoreinfo}{}
\renewcommand{\reportbodypages}{}
\urlstyle{rm}
\renewcommand{\reportwebsiteroot}{https://doi.org/10.5066/}
\renewcommand{\reportwebsiteremainder}{P1IJAXDZ}
\ifdef{\usgsissn}{\renewcommand{\usgsissn}{}}{}
\renewcommand{\theconventions}{}
\definecolor{coverbar}{RGB}{0, 47, 87}
\renewcommand{\bannercolor}{\color{coverbar}}

\renewcommand{\reportrefname}{References Cited}

\newcommand{\programname}{MODFLOW 6}

\def\tsc#1#2{\csdef{#1}{#2\xspace}}
\tsc{mff}{MODFLOW}
\tsc{mf}{MODFLOW~6}
\tsc{mfpar}{(MODFLOW~6)}

\usepackage{placeins}
\usepackage{float}
\floatstyle{plain}
\newfloat{exampleinput}{H}{exi}
\floatname{exampleinput}{}

\begin{document}

\ifdef{\makefrontcoveralt}{\makefrontcoveralt}{\makefrontcover}
\ifdef{\makefrontmatterabv}{\makefrontmatterabv}{\makefrontmatter}

\onecolumn
\pagestyle{body}
\RaggedRight
\hbadness=10000
\setlength{\parindent}{1.5pc}

\SECTION{Introduction}
% Introduction: API loop, addressing convention, and the variable taxonomy.

The MODFLOW 6 Application Programming Interface (API) provides access to the
program's internal memory while a simulation is running. The API implements the
Basic Model Interface (BMI) and an eXtended Model Interface (XMI) that exposes
the nonlinear solution loop. The API is distributed
as a shared library (\texttt{libmf6}) and is most commonly driven from Python
through the \texttt{modflowapi} and \texttt{xmipy} packages.

Every variable stored in the MODFLOW 6 memory manager is exposed through the
API. A variable is referenced by its \emph{memory address}, which has the form

\begin{quote}
\texttt{MODELNAME/PACKAGENAME/VARNAME}
\end{quote}

\noindent where all components are uppercase (for example,
\texttt{MYMODEL/NPF/CONDSAT} or \texttt{MYMODEL/STO/SS}).

\subsection{Reading and Writing Variables: Pointers versus Copies}

A variable's value is accessed through one of three routines, which differ in
both cost and semantics. Choosing the right one matters for performance in a
time loop and for correctness when a value is written.

\begin{description}
\item[\texttt{get\_value\_ptr}] returns an array that is a direct view of
MODFLOW 6's memory; no data are copied. Reading it always reflects the current
value, and assigning into it in place writes directly into MODFLOW 6's memory.
Because nothing is copied, a pointer can be obtained once (for example, just
after \texttt{initialize}) and reused for the entire simulation, which makes it
the most efficient option for a variable accessed every time step or every
iteration. The view shares MODFLOW 6's exact data type and shape, so it must be
written with the matching type and dimensions. A pointer is valid only while the
underlying array remains allocated in place; if MODFLOW 6 reallocates the
variable the pointer becomes stale and must be obtained again. Writing through a
pointer bypasses the memory-set handler described below, and pointers are not
available for string variables.

\item[\texttt{get\_value}] returns a \emph{copy} of the value. The copy is
independent of MODFLOW 6's memory: it is a snapshot that does not change when
MODFLOW 6 later updates the variable, and it cannot be invalidated by
reallocation. Each call allocates and fills a new array, costing memory and time
proportional to the variable's size, so it is best for occasional reads or when
an unchanging snapshot is wanted (for example, to compare a value before and
after a solve).

\item[\texttt{set\_value}] copies a supplied array \emph{into} MODFLOW 6's
memory. Like \texttt{get\_value} it copies element by element and pays that cost
on every call. Unlike a write through a pointer, \texttt{set\_value} invokes
MODFLOW 6's memory-set handler, so it is the appropriate choice if a variable
ever triggers a recomputation when it is set externally. No standard package
registers such a handler at present, but the mechanism exists and
\texttt{set\_value} is the forward-compatible setter.
\end{description}

\noindent In short: obtain a pointer with \texttt{get\_value\_ptr} and reuse it
when repeatedly reading or writing the same variable inside the time loop; use
\texttt{get\_value} when an independent snapshot is wanted; and use
\texttt{set\_value} for a single copy-in or when honoring a memory-set handler
matters. The pseudo-code and example in this report obtain a pointer once and
write through it, which is the common pattern for an API time loop. The named
routines follow the \texttt{xmipy} and \texttt{modflowapi} interfaces; the
\texttt{modflowapi} package additionally offers higher-level helpers, but the
cost and semantics described here are the same.

\subsection{The Simulation Loop and When Variables Are Written}

Whether a modification ``takes'' depends entirely on \emph{when} MODFLOW 6 next
writes the variable relative to when the user sets it. A MODFLOW 6 simulation
advances through a fixed sequence of phases, and each package writes its
memory-manager variables at well-defined points in that sequence:

\begin{itemize}
\item \textbf{Allocate and Read (AR)} --- executed once, after input is read.
Many derived quantities (for example, saturated conductance in the Node
Property Flow Package) are computed here from the raw input arrays and then used
unchanged for the rest of the simulation.
\item \textbf{Read and Prepare (RP)} --- executed at the start of each stress
period. List-based and array-based stress data (for example, well pumping rates)
are re-read from the input context here, overwriting any externally set values.
\item \textbf{Advance (AD)} --- executed at the start of each time step.
Time-varying inputs (for example, time-varying hydraulic conductivity, TVK) and
the quantities derived from them are refreshed here when those options are
active.
\item \textbf{Formulate (CF/FC/FN)} --- executed at every outer (Picard or
Newton) iteration. Variables read directly here (for example, specific storage
and specific yield in the Storage Package) reflect whatever value is currently
in memory at each iteration.
\item \textbf{Calculate Flows (CQ)} and output --- executed at the end of each
time step.
\end{itemize}

\noindent When using \texttt{modflowapi}, these phases correspond to callback
hooks (such as the stress-period and time-step callbacks); a value should be set
in the callback that fires \emph{after} MODFLOW 6 last writes the variable and
\emph{before} the variable is next used. As a general rule, values are set after
\texttt{prepare\_time\_step} and before the solve.

\subsection{Variable Modifiability Classes}

Each variable documented in this report is assigned one of the following
classes. The class can depend on which package options are active, as noted in
the per-variable entries.

\begin{description}
\item[Live input] An input array that is read once (or re-read each stress
period) and then used directly in the formulate step. Setting it through the API
changes the simulation, and the change persists until MODFLOW 6 next reads that
input. Example: \texttt{STO/SS}, \texttt{STO/SY}.

\item[Derived variable] A quantity computed from input during AR and then used
directly thereafter. This is the variable to set in order to change the
associated behavior; setting the raw input that produced it has no effect.
Example: \texttt{NPF/CONDSAT} (set this to change intercell conductance).

\item[Inert input] An input that is used only to compute a derived variable
during AR. Setting it after AR has no effect, because the derived quantity has
already been computed. To achieve the intended change, set the corresponding
derived variable instead. Example: \texttt{NPF/K} (when time-varying
conductivity is not active).

\item[Period-reset] A variable that is re-read from the input context at the
start of every stress period. It can be set through the API, but the value is
overwritten at the next Read and Prepare. Example: well pumping rates in
\texttt{WEL/BOUND}.

\item[Overwritten] A variable that MODFLOW 6 recomputes every time step (or
every iteration). Setting it has no lasting effect. This class includes derived
knobs whose source inputs are refreshed each time step when an option such as
TVK, time-varying storage (TVS), the Viscosity Package (VSC), or a time series
is active.

\item[State] Solution state owned by the solver or model, such as the dependent
variable (head). These can be set through the API but require care, because they
participate directly in the nonlinear solution.

\item[Structural / read-only] Grid geometry, cell connectivity, and array sizes.
These define invariants relied upon throughout the program and must be treated
as read-only. Example: \texttt{DIS/NODES}, connection arrays.
\end{description}

\subsection{How to Read the Package Tables}

Each package section lists the memory-manager variables that an API user is
likely to want to modify. For each variable the table gives its name (the
\texttt{VARNAME} component of the memory address), its modifiability class, and
guidance describing what to set, when to set it, any options that change the
classification, and the variable to set instead when the listed variable is
inert or overwritten. Variables that are purely structural or that are internal
bookkeeping are generally omitted unless they are a common source of confusion.


\SECTION{Groundwater Flow (GWF) Model Internal Flow Packages}
\subsection{Node Property Flow (NPF) Package}

The Node Property Flow (NPF) Package computes intercell conductances from the
hydraulic conductivity input and cell geometry. The single most common API
modification for this package is changing intercell conductance, which is done
by setting the derived \texttt{CONDSAT} array rather than the hydraulic
conductivity input. Under the standard (non-XT3D) formulation, the hydraulic
conductivity arrays are read only during Allocate and Read, where they are used
to compute \texttt{CONDSAT}; changing them after that point has no effect.

\begin{longtable}{p{0.15\textwidth} p{0.15\textwidth} p{0.60\textwidth}}
\caption{Memory-manager variables in the NPF Package that are candidates for
modification through the API. The memory address of each variable is
\texttt{MODELNAME/NPF/VARNAME}.} \label{table:api-gwf-npf} \\

\hline
\textbf{Variable} & \textbf{Class} & \textbf{Guidance} \\
\hline
\endfirsthead

\hline
\textbf{Variable} & \textbf{Class} & \textbf{Guidance} \\
\hline
\endhead

\texttt{CONDSAT} & Derived variable &
Saturated intercell conductance, dimensioned by the number of symmetric
connections (NJAS). Computed once during Allocate and Read from the hydraulic
conductivity arrays, cell geometry, and \texttt{ICELLTYPE}, then used directly
in every formulate step. Set \texttt{CONDSAT} to change intercell conductance at
run time. \emph{Caveats:} the array is allocated with size zero when the XT3D
option is active (conductance is then built by XT3D from the conductivity
arrays). It is recomputed every time step when time-varying hydraulic
conductivity (TVK) or the Viscosity (VSC) Package is active, which overwrites any
externally set values. \\

\texttt{K11} & Inert input &
Hydraulic conductivity along the major axis, dimensioned by the number of cells
(NODES). Used only to compute \texttt{CONDSAT} during Allocate and Read; setting
it afterward has no effect under the standard formulation---set \texttt{CONDSAT}
instead. \emph{Exceptions:} when XT3D is active the conductance is built from the
conductivity arrays during formulation, so \texttt{K11} becomes influential; when
TVK is active \texttt{K11} is re-read each time step and overwrites API changes
(class becomes Overwritten). \\

\texttt{K22} & Inert input &
Hydraulic conductivity along the second axis, dimensioned by NODES (present when
the \texttt{K22} input is specified). Same guidance and exceptions as
\texttt{K11}. \\

\texttt{K33} & Inert input &
Hydraulic conductivity along the third (commonly vertical) axis, dimensioned by
NODES (present when the \texttt{K33} input is specified). Same guidance and
exceptions as \texttt{K11}. \\

\texttt{ANGLE1}, \texttt{ANGLE2}, \texttt{ANGLE3} & Inert input &
Rotation angles of the hydraulic conductivity ellipsoid, dimensioned by NODES
(present only when specified). Used to compute \texttt{CONDSAT} (or the XT3D
coefficients) during Allocate and Read; set \texttt{CONDSAT} instead under the
standard formulation. \\

\texttt{SAT} & Overwritten &
Saturated fraction of the cell (0 to 1), dimensioned by NODES. Recomputed from
the current head at every formulate step for convertible cells, so values set
through the API are overwritten. To influence saturation, change the head
(\texttt{X}) or the cell geometry instead. \\

\texttt{ICELLTYPE} & Structural / read-only &
Cell convertibility flag (0 for confined, nonzero for convertible), dimensioned
by NODES. Read during Allocate and Read to compute \texttt{CONDSAT} and consulted
during formulation. Changing it at run time is not supported and can produce
inconsistent conductance and storage; treat as read-only. \\

\hline
\end{longtable}

\subsection{Storage (STO) Package}

The Storage (STO) Package contributes the transient storage terms to the flow
equation. Unlike the conductance in the Node Property Flow Package, the storage
capacities are not precomputed: the specific storage and specific yield arrays
are read directly during every formulate step (through the \texttt{SsCapacity}
and \texttt{SyCapacity} functions). Modifying \texttt{SS} or \texttt{SY} through
the API therefore changes the storage contribution immediately, with no derived
array to set instead.

\begin{longtable}{p{0.15\textwidth} p{0.15\textwidth} p{0.60\textwidth}}
\caption{Memory-manager variables in the STO Package that are candidates for
modification through the API. The memory address of each variable is
\texttt{MODELNAME/STO/VARNAME}.} \label{table:api-gwf-sto} \\

\hline
\textbf{Variable} & \textbf{Class} & \textbf{Guidance} \\
\hline
\endfirsthead

\hline
\textbf{Variable} & \textbf{Class} & \textbf{Guidance} \\
\hline
\endhead

\texttt{SS} & Live input &
Specific storage, dimensioned by the number of cells (NODES). Read directly at
every formulate step to compute the confined storage capacity, so changes made
through the API take effect immediately and persist. (Interpreted as a storage
coefficient rather than a specific storage when the \texttt{STORAGECOEFFICIENT}
option is active; the value set should match that interpretation.)
\emph{Caveat:} when time-varying storage (TVS) is active, \texttt{SS} is updated
each time step and overwrites externally set values. \\

\texttt{SY} & Live input &
Specific yield, dimensioned by NODES. Read directly at every formulate step to
compute the unconfined (water-table) storage capacity for convertible cells, so
changes made through the API take effect immediately and persist. Only used for
cells marked convertible by \texttt{ICONVERT}. \emph{Caveat:} when time-varying
storage (TVS) is active, \texttt{SY} is updated each time step and overwrites
externally set values. \\

\texttt{ICONVERT} & Structural / read-only &
Storage convertibility flag (0 for confined storage, nonzero for
convertible/water-table storage), dimensioned by NODES. Read every formulate step
to select the storage formulation. Changing it at run time switches a cell's
storage behavior mid-simulation and is not supported; treat as read-only. \\

\texttt{STRGSS}, \texttt{STRGSY} & Overwritten &
Cell-by-cell confined and unconfined storage flow rates, dimensioned by NODES.
These are outputs, recomputed every time step from the heads and the storage
capacities; values set through the API are overwritten. To change the storage
contribution, set \texttt{SS} or \texttt{SY} instead. \\

\hline
\end{longtable}

\subsection{Horizontal Flow Barrier (HFB) Package}

The Horizontal Flow Barrier (HFB) Package does not add a boundary term; instead it
represents thin, low-permeability barriers by \emph{modifying the intercell
conductance} between pairs of cells. It does this by editing the Node Property
Flow Package's \texttt{CONDSAT} array directly. At every Read and Prepare, the
package restores the affected connections to their pre-barrier conductance,
re-reads the barrier input, and re-applies the barrier to \texttt{CONDSAT}.

A barrier is therefore changed at run time the same way as any other intercell
conductance: by setting the appropriate entries of \texttt{NPF/CONDSAT} (the
symmetric connection index of barrier \texttt{i} is \texttt{IDXLOC[i]}, used with
the model's \texttt{JAS} array to index \texttt{CONDSAT}). The barrier's input
hydraulic characteristic, \texttt{HYDCHR}, is \emph{not} the run-time knob: it is
re-read from input and consumed at Read and Prepare (and, for the Newton
formulation or confined cells, baked into \texttt{CONDSAT} there), so a value set
through the API after \texttt{prepare\_time\_step} has no effect, except when the
XT3D option is active, where \texttt{HYDCHR} is read during formulation.

\begin{longtable}{p{0.15\textwidth} p{0.15\textwidth} p{0.60\textwidth}}
\caption{Memory-manager variables in the HFB Package that are candidates for
modification through the API. The memory address of each variable is
\texttt{MODELNAME/HFB/VARNAME}; the barrier conductance itself is changed through
\texttt{MODELNAME/NPF/CONDSAT}.} \label{table:api-gwf-hfb} \\

\hline
\textbf{Variable} & \textbf{Class} & \textbf{Guidance} \\
\hline
\endfirsthead

\hline
\textbf{Variable} & \textbf{Class} & \textbf{Guidance} \\
\hline
\endhead

\texttt{HYDCHR} & Inert input &
Hydraulic characteristic of each barrier, dimensioned by \texttt{MAXHFB}. Re-read
from input and applied to \texttt{CONDSAT} during Read and Prepare; setting it
through the API has no effect afterward (except under XT3D, where it is read
during formulation). To change a barrier at run time, set \texttt{NPF/CONDSAT}
instead. \\

\texttt{NODEN}, \texttt{NODEM} & Period-reset &
The two cells connected by each barrier, as one-based reduced node numbers,
dimensioned by \texttt{MAXHFB}. Re-read each stress period. Relocating a barrier
also requires the connection index \texttt{IDXLOC} to be made consistent, which
happens only at Read and Prepare, so moving barriers through the API is not
generally supported. \\

\texttt{NHFB} & Period-reset &
Number of active barriers in the current stress period. \\

\texttt{MAXHFB} & Structural / read-only &
Maximum number of barriers; the allocated size of the barrier arrays. Fixed when
the package is allocated and cannot be changed through the API. \\

\texttt{IDXLOC}, \texttt{CSATSAV}, \texttt{CONDSAV} & Structural / read-only &
Internal bookkeeping: the symmetric connection position of each barrier
(\texttt{IDXLOC}) and the saved conductances used to reset and re-apply the
barrier (\texttt{CSATSAV}, \texttt{CONDSAV}). These should not be set. \\

\hline
\end{longtable}

\noindent Note the interaction with \texttt{CONDSAT}: because HFB restores and
re-applies its conductance modification at every Read and Prepare, a
\texttt{CONDSAT} value set through the API on a connection that is spanned by a
barrier is overwritten at the start of each stress period. On those connections
\texttt{CONDSAT} therefore behaves as a Period-reset variable when HFB is active
and must be re-applied each step, rather than set once as described in the
``Modifying Conductance and Storage'' example.


\SECTION{Groundwater Flow (GWF) Model Stress Packages}
The list-based stress packages (Well, River, General-Head Boundary, and Drain,
among others) share a common structure, so the same guidance applies to all of
them. The per-boundary input values---pumping rates, boundary heads,
conductances, and elevations---are stored in named arrays that are re-read from
the input context at the start of each stress period during Read and Prepare.
These variables are therefore in the \textbf{Period-reset} class: a value set
through the API once before the time loop is overwritten at the next Read and
Prepare and has no effect, whereas a value set every time step (after
\texttt{prepare\_time\_step}) takes effect. The memory address of each variable
is \texttt{MODELNAME/PACKAGENAME/VARNAME}, where \texttt{PACKAGENAME} is the
user-assigned package name (not necessarily the package type).

\subsection{Variables Common to the Stress Packages}

In addition to its value variables, each package's table below includes the
per-boundary variables that are common to all of these packages and are most
often modified through the API: the auxiliary values (\texttt{AUXVAR}), the cell
of each boundary (\texttt{NODELIST}), and the number of active boundaries
(\texttt{NBOUND}). The following variables are also common to these packages but
are not listed in the per-package tables, because they are either fixed or are
recomputed by the program rather than set by the user:

\begin{description}
\item[\texttt{MAXBOUND}] The maximum number of boundaries the package can hold.
This is the allocated size of \texttt{NODELIST}, the value arrays, and
\texttt{AUXVAR}; it is fixed when the package is allocated and is in the
\textbf{Structural / read-only} class. It cannot be increased through the API.
\item[\texttt{HCOF}, \texttt{RHS}] The head-coefficient and right-hand-side
contributions of each boundary, dimensioned by \texttt{MAXBOUND}. These are
recomputed from the boundary input at every formulate step, so they are in the
\textbf{Overwritten} class. They are the mechanism by which a custom API package
injects boundary terms, but they are not the knob for a standard package---set
the boundary input variables instead.
\item[\texttt{SIMVALS}] The simulated flow for each boundary, dimensioned by
\texttt{MAXBOUND}. This is an output, recomputed every time step
(\textbf{Overwritten}).
\end{description}

\subsection{Changing Boundary Locations}

The cells a stress package acts on are not fixed: a boundary can be moved to a
different cell, and boundaries can be activated or deactivated, by editing
\texttt{NODELIST} and \texttt{NBOUND} through the API. Because these variables
are Period-reset, the same timing rule as for the value arrays applies---set them
every time step after \texttt{prepare\_time\_step}, so the change is not
overwritten by Read and Prepare.

\subsubsection{Node numbers in \texttt{NODELIST}}

The entries of \texttt{NODELIST} are one-dimensional node numbers: each is a
single flat index identifying a cell, regardless of the model's discretization
type. The cell identifier supplied in the input---layer-row-column
\texttt{(k, i, j)} for a structured (DIS) grid, layer-cell \texttt{(k, n)} for a
layered unstructured (DISV) grid, or node \texttt{(n)} for a fully unstructured
(DISU) grid---is converted to this single node number, which is what
\texttt{NODELIST} stores; it is not stored in the multi-component form used in the
input.

These node numbers are \emph{reduced} node numbers whenever any cell in the model
has an \texttt{IDOMAIN} value less than one, because cells removed from the active
grid are excluded from the numbering. When no cells are removed, the reduced node
numbers are identical to the user node numbers of the original discretization.

When reduced numbering is in effect, a \texttt{NODELIST} value cannot be
interpreted directly as a position in the original discretization; it must first
be mapped from a reduced node number to a user node number. Use the mapping arrays
\texttt{MODELNAME/DIS/NODEUSER} (reduced node number to user node number) and
\texttt{MODELNAME/DIS/NODEREDUCED} (user node number to reduced node number, zero
for removed cells), substituting \texttt{DISV} or \texttt{DISU} for \texttt{DIS}
according to the grid type. When the model has no removed cells these arrays are
allocated with a length of one, which itself signals that reduced and user node
numbers coincide. The user node number can then be converted to the
\texttt{(k, i, j)}, \texttt{(k, n)}, or \texttt{(n)} cell identifier of the
original grid from the grid dimensions; the same mapping is used in reverse to set
\texttt{NODELIST} from a known cell.

\begin{description}
\item[Moving a boundary] Set \texttt{NODELIST[i]} to the reduced node number of
the target cell. The boundary's value arrays (for example \texttt{COND} and
\texttt{BHEAD}) keep their entry index, so they continue to apply to boundary
\texttt{i} at its new cell.
\item[Activating a boundary] Increase \texttt{NBOUND} (up to \texttt{MAXBOUND})
and fill every array for the new index: \texttt{NODELIST}, the value arrays, and
any \texttt{AUXVAR}. Entries beyond the original \texttt{NBOUND} hold
unspecified memory until they are set, so all of them must be assigned.
\item[Deactivating a boundary] Decrease \texttt{NBOUND}, or move the unwanted
boundary to an inactive cell; boundaries whose cell is inactive contribute
nothing.
\end{description}

\noindent Several constraints and cautions apply. \texttt{MAXBOUND} is the hard
upper limit on the number of boundaries and cannot be grown through the API, so a
model that will gain boundaries at run time must be given a large enough
\texttt{MAXBOUND} in its input (for example, by specifying placeholder
boundaries). The target cell of a moved or activated boundary must be active.
The matrix structure is unaffected, because these packages contribute only to the
diagonal of the boundary cell, which always exists. Finally, anything keyed to a
boundary's position---auxiliary variables, boundary names, observations, and
Mover (MVR) connections---is indexed by the boundary's position in the list and
does not follow the boundary to a new cell, so relocating boundaries can
invalidate those associations.

\subsection{Well (WEL) Package}

The Well Package applies a specified volumetric flow rate to each listed cell.
The single per-boundary input is the well rate.

\begin{longtable}{p{0.15\textwidth} p{0.15\textwidth} p{0.60\textwidth}}
\caption{Memory-manager variables in the WEL Package that are candidates for
modification through the API. See also the variables common to all stress
packages described above.} \label{table:api-gwf-wel} \\

\hline
\textbf{Variable} & \textbf{Class} & \textbf{Guidance} \\
\hline
\endfirsthead

\hline
\textbf{Variable} & \textbf{Class} & \textbf{Guidance} \\
\hline
\endhead

\texttt{Q} & Period-reset &
Volumetric well rate for each well, dimensioned by \texttt{MAXBOUND} (positive
for injection, negative for extraction). Read from input at the start of each
stress period. Set it every time step (after \texttt{prepare\_time\_step}) to
change pumping at run time; a value set once before the time loop is overwritten
at the next Read and Prepare. \emph{Caveats:} when a well rate is supplied by a
time series, it is refreshed each time step and overwrites API changes (class
becomes Overwritten). When \texttt{AUTO\_FLOW\_REDUCE} is active, \texttt{Q} is
the requested rate; the rate actually applied to an extraction well is reduced as
the cell head approaches the cell bottom. \\

\texttt{AUXVAR} & Period-reset &
Auxiliary variable values, dimensioned by the number of auxiliary variables and
\texttt{MAXBOUND} (present when \texttt{AUXILIARY} variables are specified).
Re-read each stress period; set after \texttt{prepare\_time\_step}. Indexed by
boundary position, so a value follows the boundary index, not the cell. \\

\texttt{NODELIST} & Period-reset &
One-based reduced node number of the cell containing each boundary, dimensioned
by \texttt{MAXBOUND}. Set it after \texttt{prepare\_time\_step} to relocate a
boundary (see ``Changing Boundary Locations''). \\

\texttt{NBOUND} & Period-reset &
Number of active boundaries in the current stress period, from 0 to
\texttt{MAXBOUND}. Increase it (and fill the new entries) to activate boundaries
or decrease it to deactivate them (see ``Changing Boundary Locations''). \\

\hline
\end{longtable}

\subsection{River (RIV) Package}

The River Package represents a head-dependent flux between the aquifer and a
surface-water feature, limited by the riverbed bottom elevation. Its three
per-boundary inputs are the river stage, the riverbed conductance, and the
riverbed bottom elevation.

\begin{longtable}{p{0.15\textwidth} p{0.15\textwidth} p{0.60\textwidth}}
\caption{Memory-manager variables in the RIV Package that are candidates for
modification through the API. See also the variables common to all stress
packages described above.} \label{table:api-gwf-riv} \\

\hline
\textbf{Variable} & \textbf{Class} & \textbf{Guidance} \\
\hline
\endfirsthead

\hline
\textbf{Variable} & \textbf{Class} & \textbf{Guidance} \\
\hline
\endhead

\texttt{STAGE} & Period-reset &
River stage (head) for each reach, dimensioned by \texttt{MAXBOUND}. Read from
input at the start of each stress period; set it every time step to change the
stage at run time. \\

\texttt{COND} & Period-reset &
Riverbed hydraulic conductance for each reach, dimensioned by \texttt{MAXBOUND}.
Read from input at the start of each stress period; set it every time step to
change the conductance at run time. \\

\texttt{RBOT} & Period-reset &
Riverbed bottom elevation for each reach, dimensioned by \texttt{MAXBOUND}. The
exchange flux is limited when the aquifer head falls below this elevation. Read
from input at the start of each stress period; set it every time step to change
it at run time. \\

\texttt{AUXVAR} & Period-reset &
Auxiliary variable values, dimensioned by the number of auxiliary variables and
\texttt{MAXBOUND} (present when \texttt{AUXILIARY} variables are specified).
Re-read each stress period; set after \texttt{prepare\_time\_step}. Indexed by
boundary position, so a value follows the boundary index, not the cell. \\

\texttt{NODELIST} & Period-reset &
One-based reduced node number of the cell containing each boundary, dimensioned
by \texttt{MAXBOUND}. Set it after \texttt{prepare\_time\_step} to relocate a
boundary (see ``Changing Boundary Locations''). \\

\texttt{NBOUND} & Period-reset &
Number of active boundaries in the current stress period, from 0 to
\texttt{MAXBOUND}. Increase it (and fill the new entries) to activate boundaries
or decrease it to deactivate them (see ``Changing Boundary Locations''). \\

\hline
\end{longtable}

\noindent As with all of these variables, a value supplied by a time series is
refreshed each time step and overwrites API changes.

\subsection{General-Head Boundary (GHB) Package}

The General-Head Boundary Package represents a head-dependent flux to a boundary
held at a specified head through a specified conductance. Its two per-boundary
inputs are the boundary head and the conductance.

\begin{longtable}{p{0.15\textwidth} p{0.15\textwidth} p{0.60\textwidth}}
\caption{Memory-manager variables in the GHB Package that are candidates for
modification through the API. See also the variables common to all stress
packages described above.} \label{table:api-gwf-ghb} \\

\hline
\textbf{Variable} & \textbf{Class} & \textbf{Guidance} \\
\hline
\endfirsthead

\hline
\textbf{Variable} & \textbf{Class} & \textbf{Guidance} \\
\hline
\endhead

\texttt{BHEAD} & Period-reset &
Boundary head for each boundary, dimensioned by \texttt{MAXBOUND}. Read from
input at the start of each stress period; set it every time step to change the
boundary head at run time. \\

\texttt{COND} & Period-reset &
Hydraulic conductance for each boundary, dimensioned by \texttt{MAXBOUND}. Read
from input at the start of each stress period; set it every time step to change
the conductance at run time. \\

\texttt{AUXVAR} & Period-reset &
Auxiliary variable values, dimensioned by the number of auxiliary variables and
\texttt{MAXBOUND} (present when \texttt{AUXILIARY} variables are specified).
Re-read each stress period; set after \texttt{prepare\_time\_step}. Indexed by
boundary position, so a value follows the boundary index, not the cell. \\

\texttt{NODELIST} & Period-reset &
One-based reduced node number of the cell containing each boundary, dimensioned
by \texttt{MAXBOUND}. Set it after \texttt{prepare\_time\_step} to relocate a
boundary (see ``Changing Boundary Locations''). \\

\texttt{NBOUND} & Period-reset &
Number of active boundaries in the current stress period, from 0 to
\texttt{MAXBOUND}. Increase it (and fill the new entries) to activate boundaries
or decrease it to deactivate them (see ``Changing Boundary Locations''). \\

\hline
\end{longtable}

\noindent As with all of these variables, a value supplied by a time series is
refreshed each time step and overwrites API changes.

\subsection{Drain (DRN) Package}

The Drain Package removes water from the aquifer at a rate proportional to the
amount by which the aquifer head exceeds a specified drain elevation, and removes
no water when the head is below that elevation. Its two per-boundary inputs are
the drain elevation and the conductance.

\begin{longtable}{p{0.15\textwidth} p{0.15\textwidth} p{0.60\textwidth}}
\caption{Memory-manager variables in the DRN Package that are candidates for
modification through the API. See also the variables common to all stress
packages described above.} \label{table:api-gwf-drn} \\

\hline
\textbf{Variable} & \textbf{Class} & \textbf{Guidance} \\
\hline
\endfirsthead

\hline
\textbf{Variable} & \textbf{Class} & \textbf{Guidance} \\
\hline
\endhead

\texttt{ELEV} & Period-reset &
Drain elevation for each drain, dimensioned by \texttt{MAXBOUND}. Read from input
at the start of each stress period; set it every time step to change the drain
elevation at run time. \\

\texttt{COND} & Period-reset &
Drain conductance for each drain, dimensioned by \texttt{MAXBOUND}. Read from
input at the start of each stress period; set it every time step to change the
conductance at run time. When a drain-depth auxiliary variable
(\texttt{AUXDEPTHNAME}) is used, the conductance applied as the head approaches
the drain elevation is scaled from \texttt{COND} (linearly under the standard
formulation, cubically under the Newton-Raphson formulation); \texttt{COND}
remains the unscaled conductance to set. \\

\texttt{AUXVAR} & Period-reset &
Auxiliary variable values, dimensioned by the number of auxiliary variables and
\texttt{MAXBOUND} (present when \texttt{AUXILIARY} variables are specified).
Re-read each stress period; set after \texttt{prepare\_time\_step}. Indexed by
boundary position, so a value follows the boundary index, not the cell. \\

\texttt{NODELIST} & Period-reset &
One-based reduced node number of the cell containing each boundary, dimensioned
by \texttt{MAXBOUND}. Set it after \texttt{prepare\_time\_step} to relocate a
boundary (see ``Changing Boundary Locations''). \\

\texttt{NBOUND} & Period-reset &
Number of active boundaries in the current stress period, from 0 to
\texttt{MAXBOUND}. Increase it (and fill the new entries) to activate boundaries
or decrease it to deactivate them (see ``Changing Boundary Locations''). \\

\hline
\end{longtable}

\noindent As with all of these variables, a value supplied by a time series is
refreshed each time step and overwrites API changes.


\SECTION{Examples}
\subsection{Modifying Conductance and Storage}

The internal flow properties---intercell conductance in the Node Property Flow
Package and storage in the Storage Package---are the simplest variables to modify
through the API, because MODFLOW 6 does not re-read them each stress period.
\texttt{NPF/CONDSAT} is computed once during Allocate and Read and then used
directly; \texttt{STO/SS} and \texttt{STO/SY} are read directly at every formulate
step. In all three cases a value set once persists for the rest of the
simulation, so it can be assigned a single time before the time loop and the
high-level \texttt{update} routine used to advance the simulation. Set
\texttt{CONDSAT} (not the hydraulic conductivity, which is inert after Allocate
and Read) to change conductance, and set \texttt{SS} or \texttt{SY} directly to
change storage.

\begin{lstlisting}[style=pseudocode]
initialize()
set MODEL/NPF/CONDSAT = ...       # change intercell conductance (not K, which is inert)
set MODEL/STO/SS = ...            # change confined (specific) storage
set MODEL/STO/SY = ...            # change water-table storage (specific yield)
while time < end_time:
    update()                     # nothing re-reads these, so the change persists
finalize()
\end{lstlisting}

\noindent A property can also be changed partway through a simulation by setting
it before the time step from which the new value should take effect; it then
persists until it is changed again.

\begin{lstlisting}[style=pseudocode]
initialize()
while time < end_time:
    if time >= change_time:
        set MODEL/NPF/CONDSAT = ...   # takes effect from this step onward
    update()
finalize()
\end{lstlisting}

\noindent The set-once approach is valid only when MODFLOW 6 does not rewrite the
variable. If time-varying hydraulic conductivity (TVK), time-varying storage
(TVS), or the Viscosity (VSC) Package is active, the corresponding variable is
refreshed every time step and must instead be set after \texttt{prepare\_time\_step}
on each step, using the pattern shown next for stress-package input. Similarly, if
the Horizontal Flow Barrier (HFB) Package is active, it overwrites \texttt{CONDSAT}
on the connections it spans at every Read and Prepare, so those entries must be
re-applied each step rather than set once.

\subsection{Updating Stress-Package Input at Run Time}

The values in a stress package are written by MODFLOW 6 inside
\texttt{prepare\_time\_step}, which performs Read and Prepare (only on the first
time step of a stress period that reads a new period-data block) and Advance
(every time step). A value supplied as a constant in the package file is written
during Read and Prepare; a value supplied by a time series is written during
Advance. An API override must therefore be applied \emph{after}
\texttt{prepare\_time\_step}; a value set before it is overwritten by whichever
of those two steps writes the variable.

In the simplest case the value can be set before the time step and the step
advanced with the high-level convenience routine \texttt{update}, which runs an
entire time step (\texttt{prepare\_time\_step}, the iteration loop, and
\texttt{finalize\_time\_step}) in a single call. This works only when MODFLOW 6
will not rewrite the value during the step. When it would---most commonly for
time-series values---the lower-level routines are used instead so the value can
be set after \texttt{prepare\_time\_step}. The pseudo-code uses the memory
address \texttt{MODEL/PKG/COND} to stand for any stress-package input variable.

\subsubsection{Simplest use case: setting a value before the time step}

A value can be set before the time step (and the step advanced with
\texttt{update}) when MODFLOW 6 will not rewrite it during that step. This is the
case only when both of the following hold: the value is not supplied by a time
series (a time series rewrites it every step during Advance), and the step is not
the first step of a stress period that reads a new period-data block (Read and
Prepare rewrites the named array then). Mid-period time steps, and stress periods
that reuse the previous period's data, satisfy both conditions.

\begin{lstlisting}[style=pseudocode]
initialize()
while time < end_time:
    # valid only when, for this step, MF6 will not overwrite the value:
    #   - the value is not driven by a time series, AND
    #   - this is not the first step of a stress period with a new PERIOD block
    if value_will_not_be_overwritten_this_step:
        set MODEL/PKG/COND = ...
    update()                     # = prepare_time_step + iteration loop + finalize_time_step
finalize()
\end{lstlisting}

\subsubsection{Standard pattern: update after \texttt{prepare\_time\_step}}

When the value may be rewritten during the step---most importantly a time-series
value, which is refreshed every step during Advance---it must be set after
\texttt{prepare\_time\_step}. Setting it every time step at that point, using the
lower-level routines, is the safe and general pattern for both package-file and
time-series values.

\begin{lstlisting}[style=pseudocode]
initialize()
while time < end_time:
    prepare_time_step(dt)        # Read and Prepare (period start) + Advance (time series)
    set MODEL/PKG/COND = ...      # safe here: after MF6 has written its own value
    prepare_solve()
    repeat:
        converged = solve()      # one outer (Picard/Newton) iteration
    until converged
    finalize_solve()
    finalize_time_step()
finalize()
\end{lstlisting}

\subsubsection{Head-dependent values within the iteration loop}

Each \texttt{solve} call performs one outer (Picard or Newton) iteration, in which
the stress packages are formulated from their input variables. Modifying a
package input before each \texttt{solve} therefore re-linearizes the boundary at
every iteration, which can be used to implement a value that depends on the
current head (a custom head-dependent boundary or nonlinear coefficient). Modify
the package input (for example \texttt{COND}); MODFLOW 6 recomputes the
corresponding \texttt{HCOF} and \texttt{RHS} terms from it on each iteration. The
outer iteration loop converges when the head and the value it depends on reach a
consistent state.

\begin{lstlisting}[style=pseudocode]
initialize()
while time < end_time:
    prepare_time_step(dt)
    prepare_solve()
    repeat:                              # outer (Picard/Newton) iteration loop
        head = get_value_ptr("MODEL/X")
        set MODEL/PKG/COND = f(head)     # modify the package input before each solve
        converged = solve()             # this iteration formulates with the new value
    until converged                     # head and value reach self-consistency
    finalize_solve()
    finalize_time_step()
finalize()
\end{lstlisting}

A complete, runnable version of the standard pattern---with one general-head
boundary whose conductance is a constant in the package file and a second whose
conductance is supplied by a time series---is provided in
\texttt{doc/mf6api/examples/update\_stress\_via\_api.py}.



\justifying
\vspace*{\fill}
\clearpage
\pagestyle{backofreport}
\makebackcover
\end{document}
